% Generated draft: requires mathematical and source review.
\documentclass[11pt]{article}
\usepackage[margin=25mm]{geometry}
\usepackage{fontspec}
\setmainfont{Latin Modern Roman}
\usepackage{amsmath,amssymb,mathtools,unicode-math}
\setmathfont{Latin Modern Math}
\usepackage{enumitem,needspace,xurl,hyperref}
\setlength{\parindent}{0pt}
\setlength{\parskip}{5pt}
\setlength{\emergencystretch}{3em}
\newcommand{\N}{\mathbb{N}}
\newcommand{\Z}{\mathbb{Z}}
\newcommand{\Q}{\mathbb{Q}}
\newcommand{\R}{\mathbb{R}}
\newcommand{\C}{\mathbb{C}}
\newcommand{\F}{\mathbb{F}}
\newcommand{\A}{\mathbb{A}}
\newcommand{\PP}{\mathbb{P}}
\newcommand{\E}{\mathbb{E}}
\newcommand{\eps}{\varepsilon}
\newcommand{\dd}{\,\mathrm{d}}
\begin{document}
\section*{UnsolvedMath ID: 30004209}
\textbf{Source problem code:} catalogue entry
\subsection{Definitions and mathematical statement}
Let \(d\) be a finite positive integer and let \(\mathbb F_d\) denote the free group on \(d\) generators. Write \(L(\mathbb F_d)\) for its group von Neumann algebra, with its canonical trace. The supplied record does not specify the intended range of \(d\).

For a conventional interpretation of the polynomial question, let \((M,\tau)\) be a tracial von Neumann algebra generated by freely independent self-adjoint semicircular elements \(S=(S_1,\ldots,S_d)\), so that \(M=W^*(S_1,\ldots,S_d)\cong L(\mathbb F_d)\). Here a centered semicircular element of variance \(\sigma^2>0\) has spectral distribution
\[
\frac{1}{2\pi\sigma^2}\sqrt{4\sigma^2-t^2}\,\mathbf 1_{[-2\sigma,2\sigma]}(t)\,\dd t.
\]
Freeness means that alternating products of centered elements from distinct members of the family of generated unital algebras have trace zero. The source says only ``semicircular \(d\)-tuple''; free independence, generation of the ambient algebra, and normalization are not explicitly stated. The identification above therefore describes the usual free semicircular interpretation, rather than an additional verified source hypothesis.

Under the usual interpretation of polynomial evaluation, take
\[
p=(p_1,\ldots,p_d),\qquad p_j\in\C\langle X_1,\ldots,X_d\rangle,
\qquad x=p_1(S),\qquad A=W^*(x).
\]
Here \(W^*(x)\) is the unital von Neumann algebra generated by \(x\) and \(x^*\). The source does not specify the polynomial class, require the polynomials to be self-adjoint, or impose conditions on the other coordinates of \(p\). If \(A\) is abelian, then \(x\) must in particular be normal.

A unital von Neumann subalgebra \(A\subseteq M\) is a maximal abelian self-adjoint subalgebra, abbreviated MASA, if it is abelian and
\[
A'\cap M=A.
\]
Use the following tracial free-product formulation of a freely complemented MASA: there exists a unital von Neumann subalgebra \(B\subseteq M\) such that
\[
W^*(A\cup B)=M
\quad\text{and}\quad
(M,\tau)\cong(A,\tau|_A)*(B,\tau|_B),
\]
where the isomorphism respects the given inclusions. Equivalently, \(A\) and \(B\) generate \(M\) and are free with respect to \(\tau\): whenever \(n\geq1\), each \(a_k\) belongs to \(A\) or \(B\), consecutive factors belong to different algebras, and \(\tau(a_k)=0\), one has \(\tau(a_1\cdots a_n)=0\). The supplied record does not define ``freely complemented'' or specify whether the complement must additionally be nontrivial or diffuse; those conventions require confirmation.

An amenable von Neumann subalgebra \(A\subseteq M\) is maximal amenable in \(M\) if every amenable von Neumann algebra \(N\) satisfying \(A\subseteq N\subseteq M\) equals \(A\). Amenability is understood in its usual von Neumann algebra sense.

The supplied problem asks for conditions on \(p\) under which \(W^*(p_1(S))\) is a freely complemented MASA in the intended ambient algebra, and asks whether this property is generic. It also asks whether most MASAs in \(L(\mathbb F_d)\) are maximal amenable and whether every maximal amenable MASA in \(L(\mathbb F_d)\) is freely complemented. For a fixed intended \(d\), the last question has the quantified form
\[
\begin{gathered}
\text{For every MASA }A\subseteq L(\mathbb F_d),\\
A\text{ maximal amenable in }L(\mathbb F_d)
\ \Longrightarrow\
A\text{ freely complemented in }L(\mathbb F_d).
\end{gathered}
\]
The terms ``generic'' and ``most'' cannot be given precise quantifiers from the supplied record: no parameter space, topology, measure, or equivalence convention is specified. Thus these portions remain formulation questions, and even the displayed implication depends on resolving the intended range of \(d\) and complement convention.
\subsection{Short English statement}
The problem asks when a polynomial in a semicircular family generates a freely complemented MASA, whether this happens generically, whether most MASAs in \(L(\mathbb F_d)\) are maximal amenable, and whether every maximal amenable MASA is freely complemented. The source leaves the polynomial class, several ambient conventions, and the meanings of ``generic'' and ``most'' unspecified. This is a statement-only draft requiring human mathematical review.
\subsection{Sources}
\begin{itemize}
\item Supplied local UnsolvedMath registry record, permanent ID \(30004209\), title ``Freely Complemented and Maximal Amenable MASAs,'' source problem code \texttt{catalogue entry}; collection: Oberwolfach Reports Open Problems. Project registry locations are \texttt{data/unsolvedmath/problems.json} and \texttt{catalog/problems.json}. The supplied record is the evidence used here; no independent registry inspection is claimed.
\item Supplied source citation: \emph{C*-Algebras (2020)}, Oberwolfach Reports, DOI \texttt{10.4171/owr/2019/37}, pp.\ 2257--2332, \url{https://doi.org/10.4171/owr/2019/37}. The background describes these reports as documenting mathematical workshops and their open-problem discussions. The record gives proposed year \(2019\) and names no individual proposer.
\item Supplied literature source: ``Oberwolfach Report 37, MASA questions,'' \url{https://doi.org/10.4171/owr/2019/37}.
\item Supplied literature source: ``Boschert et al., freely complemented MASAs,'' \url{https://arxiv.org/abs/2411.05136}. No fuller author list or bibliographic title is supplied.
\item Literature assessment dated exactly \texttt{2026-08-22}: the supplied curation labels the problem open and states that the question for general polynomial images \(p_1(S)\) remains open; known generator, radial, perturbative, and reassembly examples do not settle genericity or all maximal-amenability questions. This is supplied dated literature triage derived from the OWR curation, not an independently verified assessment. Its time-sensitive status and cited claims require independent verification.
\item The supplied registry labels the statement \texttt{corrected\_verified}, explaining that mathematical content was preserved while notation, typography, duplicate clauses, or omitted definitions were repaired for self-contained publication. This records the supplied editorial qualification and does not establish independent mathematical verification of this draft.
\end{itemize}
\end{document}
